% witness: tabularray/tabularray (tabularray.tex:2859-2870 `\begin{tblr}[evaluate=\makeEmptyTable]{hlines,vlines}`,
%   :1670-1679 `\SetTblrOuter{expand=…}` + `[expand+=…]`, :2833 `evaluate=\myFunOne`, :2873 `evaluate=all`)
% oracle:  pdflatex 0 errors; tabularray.sty:3567-3575 `\__tblr_modify_table_body:` evaluates the functional library's
%   functions (:8289-8306, `all` :8275-8277) and expands each `expand=` macro once (:3579-3591) before splitting the body
% engines: rust 59l = `\makeTable{3}{4}` (rows that start with an empty cell, as the manual's `\makeEmptyTable`) ran
%   inside the first cell: "Stray alignment" (1 error; the manual's own empty table, its RED
%   tblr_all_empty_table_keeps_its_rules); the expand and evaluate=all tables were already right (the kernel tabular
%   expands macros) and stay as controls
% expect:  0 errors; tables `(empty) y y y` x 3, AA/20/50/DD x 3, Row1/All, and a/20/50 x 2 — the last one tells the
%   expansion apart from the kernel's own (the body redefines \rowa/\rowb as it runs: 59l printed X/Y, P/Q; pdflatex
%   expands first: 20/30, 50/60), and the default `expand=\rowa` merging with `expand+=\rowb` (59n review)
% status:  RED (59l) -> GREEN (59n)
\documentclass{article}
\usepackage{tabularray}
\UseTblrLibrary{functional}
\newcommand*\tblrrowa{20 & 30 & 40 \\}
\newcommand*\tblrrowb{50 & 60 & 70 \\}
\newcommand*\tblrbody{\tblrrowa \tblrrowb}
\newcommand*\rowa{20 & 30 \\}
\newcommand*\rowb{50 & 60 \\}
\begin{document}
\IgnoreSpacesOn
\prgNewFunction \makeTable {mm}  {
  \tlSet \lTmpaTl {\intReplicate {\intEval{#2-1}} {&y}}
  \tlPutRight \lTmpaTl {\\}
  \intReplicate {#1} {\tlUse \lTmpaTl}
}
\prgNewFunction \myFunOne {m} {
  \prgReturn {#1 & #1 \\}
}
\IgnoreSpacesOff
\begin{tblr}[evaluate=\makeTable]{hlines,vlines}
  \makeTable{3}{4}
\end{tblr}

\SetTblrOuter{expand=\tblrbody\tblrrowa}
\begin{tblr}[expand+=\tblrrowb]{ccc}
 AA & BB & CC \\
 \tblrbody
 DD & EE & FF \\
\end{tblr}

\begin{tblr}[evaluate=all]{hlines}
  Row1 & 1 \\
  \myFunOne{All}
\end{tblr}

\SetTblrOuter{expand=\rowa}
\begin{tblr}[expand+=\rowb]{cc}
 \expandafter\gdef\csname rowa\endcsname{X & Y \\}\expandafter\gdef\csname rowb\endcsname{P & Q \\}a & b \\
 \rowa
 \rowb
\end{tblr}
\end{document}
